\chapter{Hadamard singularities and distributional tests}\label{app:C}
At separated spacelike points the free vacuum two-point function can look
like an ordinary function of distance. At null separation its prescription
matters, and at coincidence it cannot be evaluated at a point. These are
different issues. The Hadamard form specifies a local singularity together
with a frequency orientation. The wavefront set tests whether one can
restrict or multiply the resulting distribution. Scaling degree answers a
third question: if an expression is already defined off coincidence, what
freedom remains in extending it? We give the tests for a real scalar in
the smooth bulk, where all three questions already arise.

\section{World function, frequency and normalization}
Work first in a smooth geodesically convex bulk neighborhood in four
spacetime dimensions with signature $(-,+,+,+)$. For the unique geodesic
$\gamma$ from $x'$ to $x$, affinely parametrized on $[0,1]$, Synge's world
function is
\[
 \sigma_g(x,x')=\tfrac12\int_0^1 g(\dot\gamma,\dot\gamma)\,ds .
\]
Thus it is \emph{half} the signed squared interval, negative for a timelike
separation. It is distinct from the antisymmetric form on CCR labels. Let
$T$ be a smooth time function increasing to the future, in units of length.
Use the boundary value, for $\epsilon\downarrow0$,
\[
 \sigma_\epsilon=\sigma_g+i\epsilon(T(x)-T(x'))+\epsilon^2/2 .
\]
For $T=t$ in flat space this is exactly
$\big(|\mathbf x-\mathbf x'|^2-(t-t'-i\epsilon)^2\big)/2$. Use the
branch of the logarithm real on the positive real axis. With $\hbar=1$ and
$\ell>0$ a length, the local singular part of the \emph{unordered}
(Wightman) two-point function is written
\begin{equation}
 H_\ell(x,x')=\frac{1}{8\pi^2}\lim_{\epsilon\downarrow0}
  \left[\frac{U(x,x')}{\sigma_\epsilon(x,x')}
    +V(x,x')\log\!\frac{\sigma_\epsilon(x,x')}{\ell^2}\right].
 \label{eq:C-hadamard}
\end{equation}
$U,V$ are geometric coefficients for the chosen Klein--Gordon operator,
not functions to be fitted to a state. For example $U$ obeys
\[
 2\nabla^a\sigma_g\nabla_a U+(\Box\sigma_g-4)U=0,\qquad U(x,x)=1 .
\]
To see the first equation, act with $\Box$ on $U/\sigma_g$ away from the
cone and use $\nabla_a\sigma_g\nabla^a\sigma_g=2\sigma_g$. The
coefficient of $\sigma_g^{-2}$ must vanish. The remaining inverse power
determines the leading logarithmic coefficient by transport; higher
coefficients can be constructed successively. On a smooth, not necessarily
analytic, metric this notation denotes a local parametrix modulo smooth
terms. It can be used to any specified differentiability order by
truncating its asymptotic expansion sufficiently late; convergence of an
infinite power series is not assumed \cite{WaldQFT,HWick}.

The flat massless benchmark is $U=1,V=0$. It can also be calculated from
positive-frequency modes: with $\tau=t-t'$, $r=|\mathbf x-\mathbf x'|$,
\begin{align*}
 W_0(x,x')&=\lim_{\epsilon\downarrow0}
 \int\frac{d^3p}{(2\pi)^3 2|\mathbf p|}
 e^{-i|\mathbf p|\tau+i\mathbf p\cdot(\mathbf x-\mathbf x')-\epsilon|\mathbf p|}\\
 &=\lim_{\epsilon\downarrow0}\frac{1}{4\pi^2 r}
       \int_0^\infty e^{-(\epsilon+i\tau)p}\sin(pr)\,dp
 =\lim_{\epsilon\downarrow0}\frac{1}{4\pi^2[r^2-(\tau-i\epsilon)^2]} .
\end{align*}
This checks both the $1/(8\pi^2)$ in \eqref{eq:C-hadamard} and the sign of
the prescription. Reversing the arguments reverses the frequency
orientation. Time ordering instead selects $W_0(x,x')$ for $\tau>0$ and
$W_0(x',x)$ for $\tau<0$; its flat contraction is
$G_F=(4\pi^2)^{-1}(r^2-\tau^2+i0)^{-1}$, as used in
Chapter~\ref{ch:13}. The Wightman function is a bisolution, whereas this
$G_F$ has an $i\delta$ source. Other spacetime dimensions have different
Hadamard powers or logarithms, not the four-dimensional prefactor printed
here.

A scalar Hadamard state has $W-H_\ell$ smooth locally (or the corresponding
finite regularity after truncation). Changing $\ell$ shifts the smooth
representative and finite local subtractions, not the leading singularity.
The parametrix is not in general an exact positive bisolution.
To specify an actual centered quasifree state the full $W$ must satisfy
the equation and the true boundary condition in both entries,
$W(x,x')-W(x',x)=iE(x,x')$, Hermiticity and positivity as in
Chapter~\ref{ch:11}. Two Hadamard two-point functions for the same ordinary
globally hyperbolic free problem differ by a smooth symmetric bisolution.
Their fluctuation predictions can therefore differ even if every
short-distance subtraction is the same.

\section{Which directions are singular?}
A distribution $u$ is smooth near a point if its localized Fourier
transform decays faster than any inverse power in every momentum
direction. The \emph{wavefront set} $\operatorname{WF}(u)$ records pairs
$(x,k)$ with $k\ne0$ where this test fails in direction $k$: no cutoff
nonzero at $x$ gives rapid decay in an open conical neighborhood of $k$.
We take Fourier phases $e^{-ik\cdot x}$ for this test. A positive-frequency
mode has oscillatory covector $k=(-|\mathbf p|,\mathbf p)$, so its
metric dual $k^\sharp$ is future directed. For a state of the ordinary
free scalar on a smooth globally hyperbolic spacetime the Hadamard
condition is equivalently the microlocal spectrum condition \cite{Rad}:
\begin{equation}
 \operatorname{WF}(W)=\{(x,k;x',-k'):(x,k)\sim(x',k'),\quad
                              k^\sharp\ \text{future null}\} .
 \label{eq:C-wavefront}
\end{equation}
Here $\sim$ means that $k,k'$ are cotangents to the same null geodesic
(with metric dual tangent to it) related by parallel transport. The
coincident pairs $x=x',\ k=k'$ on a null direction are included. The minus
on the second covector comes from the two arguments, not a past-directed
first frequency. A mass does not move these singular directions off the
light cone: the principal part of the operator is unchanged. At a genuine
timelike boundary reflected singularities may need to be included in a
global criterion; \eqref{eq:C-wavefront} is the standard smooth-bulk
result, not an assertion about arbitrary characteristic or corner problems.

In the ordinary bulk problem this gives a short proof of the smoothness of
state differences. Let $D=W_1-W_2$ for two Hadamard states with the same
commutator. The wavefront set of the difference is contained in the union
of their wavefront sets, so its first covector must be future null. But
$D(x,x')=D(x',x)$ by the commutator. Exchanging the arguments in
\eqref{eq:C-wavefront} would make that same first covector past null.
There are no such nonzero covectors; hence $D$ is smooth. The proof shows
that smoothness fixes neither the size nor the sign of a difference in
fluctuations.

\section{Test before restricting or multiplying}
Let $\iota:S\to M$ be a smooth embedding. Its conormal consists of the
nonzero covectors at $\iota(s)$ which annihilate every tangent to the
image. The standard pullback $\iota^*u$ is defined by the wavefront
criterion if
\[
 \operatorname{WF}(u)\ \cap\ N^*\iota=\varnothing .
\]
For a kernel of two variables apply this criterion to the actual map, for
example $\iota\times\iota:S^2\to M^2$. Both cotangents must annihilate
their respective tangent spaces to give an obstruction. On a spacelike
Cauchy surface the conormals are timelike, whereas both covectors in
\eqref{eq:C-wavefront} are null. The restriction of $W$ to two Cauchy
arguments is therefore a distribution on the product of the surface with
itself. Spacetime derivatives of $W$ do not enlarge its wavefront set, so
normal derivatives can be taken before this pullback. This is how
field--momentum as well as field--field covariances can be given by Cauchy
data. It does not define the value $W(x,x)$.

Now take the null plane $u=t-x^1=0$ and two points on the same generator,
with the same transverse coordinates. A future null covector at the first
point is $-a\,du$, $a>0$; parallel transport gives the same covector at
the second point. The pair $(-a\,du,+a\,du)$ lies both in the wavefront
set of $W$ and in the conormal to the product of null planes. The usual
pullback theorem no longer allows us to restrict both fields in this
manner. It does not prove that no characteristic observable whatsoever
exists. Smeared fluxes, specially differentiated data or other trace
constructions require their own domains and checks. The classical
existence of a smooth trace for a particular solution is not the quantum
distributional check.

Multiplication is pullback to the diagonal of a tensor product. A
sufficient condition for $uv$ is that there be no $(x,k)\in
\operatorname{WF}(u)$ with $(x,-k)\in\operatorname{WF}(v)$. Smooth
switching functions always pass this test. For a sharp step across a smooth
face $s=0$, $ds\ne0$, the wavefront set of $\Theta(s)$ contains both
normal directions $a\,ds$, $a\ne0$. Hence, for example,
$\Theta(s)\delta(s)$ fails this test: assigning a fraction of a surface
contact to one side by writing $\Theta(0)=1/2$ is not a distributional
proof. The problem can be seen by smoothing the step and shifting its
transition by a distance of the order of its width. Its product with the
delta depends on the value selected at zero even when all these smooth
steps have the same limit. Conversely, a step in only the first argument
of a bidistribution has zero cotangent in the other slot in this product
test. One must check the actual distributions, not declare every time
ordering or every step inadmissible.

This criterion also explains the puncture in the fish calculation. Two
factors of $W$ can multiply as distributions in the \emph{pair} of
variables: their nonzero first covectors are both future directed. Away
from the diagonal $G_F$ selects $W$ or its transpose locally near each
singular pair, so its square there likewise passes the test. But on the
diagonal $\Box_x G_F=i\delta(x,y)$ in the flat example. A differential
operator cannot add directions to a wavefront set, and the delta has all
diagonal conormals $(k,-k)$ there, with both signs of $k$. The square of
$G_F$ at coincidence therefore fails the product test. Extending an
already defined off-diagonal distribution, not assigning a value to a
step at zero, is the operation that remains.

\section{How much freedom does extension leave?}
Suppose a scalar distribution $t^\circ$ is given for $z\ne0$ in
$\mathbb R^k$. For $\psi$ a test function set
$\psi_a(z)=a^{-k}\psi(z/a)$, $t_a(\psi)=t^\circ(\psi_a)$ on the
punctured space. The scaling degree is
\[
 \operatorname{sd}(t^\circ)=
 \inf\{d:a^d t_a\longrightarrow0\ \text{as }a\downarrow0
                     \text{ in distributions}\} .
\]
A homogeneous singularity $|z|^{-b}$ has scaling degree $b$ off zero;
logarithms do not change the infimum. The bulk extension theorem for finite
scaling degree $\rho$ says that an extension of degree $\rho$ exists and
is unique at that degree if $\rho<k$ \cite{BF}. If $\rho\ge k$, two such
extensions may differ by
\[
 \sum_{|\alpha|\le\lfloor\rho-k\rfloor} c_\alpha\,
                                \partial_z^\alpha\delta^{(k)}(z) .
\]
Indeed their difference is supported at zero; distributions so supported
are finite sums of delta derivatives, and each derivative has degree
$k+|\alpha|$. To see why the bound also suggests a construction, subtract
from the test function its Taylor polynomial through degree
$\lfloor\rho-k\rfloor$, multiplied by a smooth cutoff equal to one near
zero. The remainder vanishes fast enough that removing a small neighborhood
of zero has a limit. Assigning the finitely many missing Taylor moments
amounts to choosing the constants $c_\alpha$. An extension is a new
prescription; it need not preserve exact homogeneity, even when it
preserves the scaling degree.

At a smooth collision stratum use local coordinates $(X,z)$ with $z=0$
on the stratum. Scale the $k$ normal coordinates $z$ only, keeping the
tangential variables $X$ as test parameters. The analogous ambiguities
are normal delta derivatives with coefficients on the stratum; in general
those coefficients are distributions subject to further requirements.
Local covariance, allowed field dimensions, unitarity, causal
factorization and Ward identities reduce the freedom in an actual field
theory. This counting works only once the distributions away from the
stratum and the requisite normal estimates are known. A null pullback that
failed the wavefront test is not automatically defined by quoting the
codimension of a coincident-point stratum. Chapter~\ref{ch:13} instead
starts with a defined product away from coincidence and extends it, which
is precisely the situation addressed by the theorem.

\begin{exercise}
For the flat massless vacuum restrict the two arguments to $t=t'=0$.
Starting from the momentum integral above, find the singularity as the
spatial separation $r$ tends to zero. Does the distribution on the pair of
Cauchy surfaces need an arbitrary spatial contact term at $r=0$ if its
scaling degree is to be preserved? Does this give a finite value for the
field squared at a point?
\end{exercise}
\noindent\emph{Hint.} The equal-time pullback passed the conormal test.
Compare $1/r^2$ with the radial volume $r^2 dr\,d\Omega$ in three relative
spatial dimensions; distinguish a locally integrable distribution from its
pointwise value at the origin.
